\documentclass[11pt]{article}
\usepackage[letterpaper,landscape,margin=1.3cm]{geometry}
\usepackage{graphicx}
\usepackage{enumitem}
\usepackage{xcolor}
\usepackage{microtype}
\usepackage[hidelinks]{hyperref}

\setlist{nosep, leftmargin=1.6em}
\setlength{\parindent}{0pt}
\setlength{\parskip}{0.25em}
\pagestyle{empty}
\newcommand{\code}[1]{\texttt{#1}}
\newcommand{\wire}[2]{\textcolor{#1}{\rule[0.25ex]{1.6em}{1.6pt}}\,\textbf{#2}}
\definecolor{wred}{RGB}{220,20,20}
\definecolor{wwhite}{RGB}{170,170,170}
\definecolor{wgreen}{RGB}{10,170,50}
\definecolor{worange}{RGB}{250,140,0}
\definecolor{wyellow}{RGB}{220,190,0}
\definecolor{wblue}{RGB}{30,90,250}
\definecolor{wgrey}{RGB}{154,160,166}
% the debugger harness, colours as delivered
\definecolor{hOrange}{HTML}{EE8A5C}
\definecolor{hPurple}{HTML}{5A2F8C}
\definecolor{hBrown}{HTML}{5B4038}
\definecolor{hBlack}{HTML}{25211F}
\definecolor{hYellow}{HTML}{E8CF2A}
\definecolor{hTeal}{HTML}{1F7A68}
\newcommand{\pic}[1]{\begin{center}\includegraphics[width=\linewidth,height=0.56\textheight,keepaspectratio]{figures/power-bench-#1.pdf}\end{center}\vspace{-0.8em}}
\newcommand{\head}[2]{{\Large\bfseries #1}\hfill{\small 2 November 2026, revision 3 \quad #2}\par\smallskip}

% the part shared by the three microcontroller boards: wired once, left wired between their sessions
\newcommand{\fixturesteps}{%
\item A \textbf{standard USB-A to USB-C cable} from the laptop into the \textbf{spare F411}'s USB-C socket.
\item \wire{wred}{5V}: spare F411 \code{5V} (the one next to \code{B9}) $\to$ M1 \code{IN+}; M1 \code{IN-} $\to$
breakout \code{VBUS}. M1's pins all sit on its one header and point down: the wires go on from underneath.
\item \wire{wwhite}{A11$\to$D$-$}, \wire{wgreen}{A12$\to$D+}, \wire{black}{G$\to$GND}: spare F411 to your breakout
(twist the two data wires).
\item Spare F411 to M1: \wire{wyellow}{B6$\to$SCL}, \wire{wblue}{B7$\to$SDA}, \wire{worange}{3V3$\to$VCC},
\wire{black}{GND$\to$GND} (the \code{GND} on its small 4-pin end header), \wire{wgrey}{5V$\to$VBS} (its \emph{other}
\code{5V}, at the USB end; sense only, the logger subtracts the shunt drop). \code{ALE} stays free.
\item Debugger to the spare F411, by its own harness colours: \wire{hOrange}{TXD$\to$A3}, \wire{hPurple}{RXD$\to$A2},
\wire{hBrown}{GND$\to$G} (the \code{G} next to \code{5V}, USB end); the debugger in the laptop's other USB port. Its red
(\code{3V3}) and light-grey (\code{5V}) wires are never connected.}
\newcommand{\fixturechecks}{%
\textbf{Steps 1--5 are the same on all three microcontroller pages:} leave them wired; between sessions only the
board, its cable and the debugger's three SWD wires change. Wire by the printed pin labels; the pictured pin positions
are approximate.}

\begin{document}

%=========================================================== overview page
\head{Power measurement wiring --- all four boards}{overview}

\textbf{One board at a time.} Each page below is one session: wire it, measure, move to the next. \textbf{Nothing is
cut, nothing is bought, and every cable is a standard USB cable.} The Pi needs no wiring at all: it reports its own
power (page 5).

\begin{center}
\begin{tabular}{@{}lllll@{}}
\textbf{page} & \textbf{board measured} & \textbf{module} & \textbf{debugger's SWD wires} & \textbf{new image onto the board}\\
\hline
2 & STM32H723 & M1, \code{R100} & to the H723 & by script, debugger stays wired\\
3 & ESP32-S3 & M1, \code{R100} & unused & over \code{COM}, then \code{COM} unplugged\\
4 & main STM32F411 & M1, \code{R100} & to the main F411 & by script, debugger stays wired\\
5 & Raspberry Pi 5 & none: its own power telemetry & unused & copied over the network, as always\\
\end{tabular}
\end{center}

\textbf{One extra board, always the same one: the spare F411.} It is the USB breakout and the logger at once, in every
session. Its USB-C socket is wired to its own header pins \code{5V}, \code{G}, \code{A11} (USB D$-$) and \code{A12}
(USB D+). It runs the logger firmware, which never switches its own USB on and leaves \code{A11}/\code{A12} idle, so the
USB lines pass straight through it to your USB-A female breakout, with the 5\,V line going through the current-sense
module. The measured board's usual cable plugs into the breakout. The logger reads the module over I$^2$C and sends the
readings to the laptop through the debugger's serial port. Its own current is drawn before the module, so it is not
part of any reading.

\textbf{The debugger stays wired for good.} Its harness wires keep their delivered colours, and they do not follow the
usual convention: \textbf{black is SWDIO and red is 3V3}. Its \wire{hOrange}{TXD}, \wire{hPurple}{RXD} and
\wire{hBrown}{GND} go to the spare F411 (the readings); its \wire{hBlack}{SWDIO}, \wire{hYellow}{SWCLK} and
\wire{hTeal}{GND} go to whichever STM32 is measured, so a script can flash one
image after another without touching a wire. \textbf{One catch:} when OpenOCD connects it switches on the chip's debug
clocks, and they stay on through a reset; only a power cycle clears them. So the flash script switches them off again
(writes 0 to \code{DBGMCU\_CR} and reads it back) and quits OpenOCD before any reading is taken. Check it once per STM32
board: the idle reading after a scripted flash must equal the idle reading after unplugging and replugging the board's
cable with the SWD wires off. If the two differ, replug the board's cable after every flash instead.

\textbf{Check once, before the first session:} with only the spare F411 plugged in and running the logger,
\code{dmesg -w} prints nothing (an \code{error -71} line means something sits on the USB data lines), its \code{5V}
pin reads about 4.7\,V (after its input diode; the module records the actual voltage), and the readings arrive on the
debugger's serial port. Then stream one utterance through the whole chain and confirm the output is bit-exact.

\textbf{Every session:} \code{IN+} towards the supply, \code{IN-} towards the board, one common ground; run the exact
image whose timing figures are in the results; two minutes' settling, then one minute each of idle, active (unpaced)
and a short paced run as the check.

{\scriptsize\color{black!60} The pictures are drawn to scale from the parts' own drawings (\code{docs/figures/},
\code{bench-parts.tex}; the same figures go into the paper). The debugger's harness is drawn in its delivered colours;
the INA226 module is drawn from a photo of the one on the bench; only the USB-A breakout is still generic. Revision 1 used both
F411s in rotation; revision 2 merges the breakout and logger jobs into the spare F411; revision 3 (same day) measures
the Pi with its own telemetry.}

%=========================================================== H723
\newpage
\head{Session 1: STM32H723}{page 2 of 5}
\pic{h723}
\begin{minipage}[t]{0.53\linewidth}
\begin{enumerate}
\fixturesteps
\end{enumerate}
\end{minipage}\hfill
\begin{minipage}[t]{0.44\linewidth}
\begin{enumerate}[start=6]
\item Debugger to the H723's SWD header (\code{GND CLK DIO 3.3V}, top edge): \wire{hBlack}{SWDIO$\to$DIO},
\wire{hYellow}{SWCLK$\to$CLK}, \wire{hTeal}{GND$\to$GND}; its \code{3.3V} pin left free.
\item The H723's usual USB-A to USB-C cable into your breakout, and into the H723. Flash the measurement image
\code{0b3cc69b} by script.
\end{enumerate}
\smallskip
{\small\fixturechecks}
\end{minipage}

%=========================================================== S3
\newpage
\head{Session 2: ESP32-S3}{page 3 of 5}
\pic{s3}
\begin{minipage}[t]{0.53\linewidth}
\begin{enumerate}
\fixturesteps
\end{enumerate}
\end{minipage}\hfill
\begin{minipage}[t]{0.44\linewidth}
\begin{enumerate}[start=6]
\item The debugger's SWD wires stay unconnected: the S3 is flashed over \code{COM}. Flash the measurement image
(\code{1f3cbc53}) over \code{COM}, then \textbf{unplug \code{COM}}: with both cables in, part of the board's current
would come in through the second one and M1 would not see it. A new image means plugging \code{COM} back in for the
flash.
\item The S3's usual USB-A to USB-C cable into your breakout, and into \textbf{the S3's port marked \code{USB}} (the left-hand one with the antenna away from you; \code{COM} is the right-hand one).
\end{enumerate}
\smallskip
{\small\fixturechecks}
\end{minipage}

%=========================================================== F411
\newpage
\head{Session 3: the main STM32F411}{page 4 of 5}
\pic{f411}
\begin{minipage}[t]{0.53\linewidth}
\begin{enumerate}
\fixturesteps
\end{enumerate}
\end{minipage}\hfill
\begin{minipage}[t]{0.44\linewidth}
\begin{enumerate}[start=6]
\item Debugger to the main F411's 4-pin SWD end header (\code{3V3 DIO SCK GND}): \wire{hBlack}{SWDIO$\to$DIO},
\wire{hYellow}{SWCLK$\to$SCK}, \wire{hTeal}{GND$\to$GND}; its \code{3V3} pin left free.
\item The main F411's usual USB-A to USB-C cable into your breakout, and into the main F411. Flash its measurement
image by script.
\end{enumerate}
\smallskip
{\small\fixturechecks{} Two F411s are on the bench in this session only because one of them is the board measured;
the spare F411 does exactly what it does in the other sessions.}
\end{minipage}

%=========================================================== Pi 5
\newpage
\head{Session 4: Raspberry Pi 5}{page 5 of 5}

\textbf{Nothing to wire.} The Pi 5's power-management chip reports the voltage and current of each of its supply rails,
\code{vcgencmd pmic\_read\_adc}. The Pi's power is the sum of voltage times current over those rails. \textbf{Decided
2 November 2026:} the Pi is measured this way rather than with a module in its supply cable, so no part has to be bought
and no session waits for one.

\textbf{What the number leaves out, and the paper says so.} The rails are measured \emph{after} the chip's own converters,
so the Pi's figure does not contain the conversion losses, or anything fed straight from the 5\,V input (USB ports, the
fan). Every microcontroller figure is taken at the 5\,V input and does contain its regulator's losses. The Pi therefore
looks better than it would on the same footing, by roughly the converters' loss. State it beside every Pi energy figure.

\begin{enumerate}
\item Official supply and cable as always; Ethernet; HDMI and USB devices unplugged.
\item Sample \code{vcgencmd pmic\_read\_adc} once a second from a script \textbf{pinned to a core the benchmark does not
use}: the sampler runs on the machine it measures, so keep it slow and off the measured core. Log the fan's PWM setting
beside it.
\item Two minutes' settling, then one minute each of idle, active (unpaced) and a short paced run, as on the other
boards; \code{vcgencmd get\_throttled} must read \code{0x0} before and after.
\end{enumerate}

{\small\textbf{Check:} sum the rails at idle and compare with a typical Pi 5 idle draw; a figure far off means a rail is
missing from the sum or a unit is wrong.}

\end{document}
